\section{Conclusions and limitations}
\label{sec:conclusions}

\subsection{What we found}
\paragraph{Analytical results.} For a stock fished with a fixed quota $H$, a strong Allee effect creates a pair of steady states, an unstable threshold and a stable upper state, which approach each other as the quota rises and merge at the critical catch $H_c=rKq_c$ (Section~\ref{sec:foundations}). As the quota approaches this limit, the gap between them, which is the stock's safety margin, shrinks like $\sqrt{q_c-q}$, so a stock fished close to the limit can be tipped into collapse by a very small shock. A fixed fraction of the stock (Model B) has an analogous critical effort $h_c=r(1-a)^2/(4a)$ (Section~\ref{sec:wp02}). The two rules must be compared in common units and with a stated calibration: with the same first-year catch the logistic stock collapsed in our example while both Allee stocks survived, but with equal fractions of their own thresholds all three survived, so no universal ranking follows.


\paragraph{Collapse and recovery.}
Reducing harvesting after a collapse does not necessarily
restore the population. Once it has fallen below the
Allee threshold, natural growth remains negative even
without harvesting. Recovery may therefore require
an intervention such as restocking the population above
the relevant recovery threshold
(Section~\ref{sec:wp02}).
Section~\ref{sec:wp03} also showed how the initial
population determines survival, how extinction times
can be calculated, and why collapse may be delayed
when constant harvesting is just above its critical rate.


\paragraph{Forward PINN.} A small tanh network with the initial condition built in and a physics-only loss reproduced the solution of all six test cases to a relative $L^2$ error between $2.9\times10^{-7}$ and $4.1\times10^{-3}$ (36 of 36 runs accurate over two sets of seeds). It was much slower than a standard solver, which reaches about $10^{-10}$ in a few milliseconds, so for a single scalar equation the PINN is not the better tool (Section~\ref{sec:wp04}).

\paragraph{Inverse PINN.} Using one shared setting for the forward and inverse problems, the PINN recovered every fitted parameter within $5\%$ in 10 of 18 clean fits, against 18 of 18 for ordinary least squares and 0 of 18 for the settings first tried (Section~\ref{sec:wp05}). Under noise, recovery was limited by the data themselves. Noise was not the reason some fits failed; the fits that did not converge came from validation-based checkpoint selection, and the remaining recovery failures were the poorly conditioned cases. The one provably impossible case, Model B with all four parameters unknown, was kept separate from these.

\subsection{Lessons about method}
\begin{itemize}
\item A small training loss is not evidence of accuracy. The early forward PINNs had a small loss at the collocation points and a large error between them.
\item A good trajectory fit does not guarantee good parameters. The inverse PINN fitted every training set to about $10^{-4}$ and still got some parameters wrong by $100\%$.
\item Check convergence against what is achievable. Comparing the training error with the noise level exposed runs that had stopped without learning.
\item A classical baseline should always be run. Here it was faster and more accurate than the PINN on every clean case.
\end{itemize}

\subsection{Limitations}
\begin{itemize}
\item All data are synthetic and generated from the same models that are fitted, so there is no model error and nothing here validates the models for a real species.
\item The models are deterministic and unstructured. Random fluctuations, age structure, delays, spatial effects and changing parameters can all matter most near a threshold.
\item The PINN studies use one parameter set per case and three seeds per setting. Differences between settings are therefore not statistically tested, and the choice of the informative-case threshold was made after seeing the results.
\item In the inverse study $K$ and the starting population are assumed known and exact, which makes the problem easier than a field setting.
\item The iteration budget of the L-BFGS phase was binding in every PINN run, and longer runs were not tried.
\item Results depend on the platform and library versions (recorded in the run manifests), and the long sweeps are recorded evidence rather than re-run for this report.
\end{itemize}

\subsection{Future work}
Natural next steps are to add random fluctuations and age structure to the model, to test other parameter values and more observations in the inverse problem, to try bounds or a different parametrisation to stop $r$ and $A$ collapsing together in the near-fold cases, and, if real data become available, to fit a real fishery with its provenance documented.